\section{Global macro: strategy, record, risk and practitioners}

\begin{qfdefinition}[Global macro: a top-down, unconstrained strategy; discretionary, systematic, high-frequency and CTA approaches; directional, relative-value and thematic trades; the four product groups and multi-asset expression of a view]\label{def:global-macro-strategy}
\qffield{Statement} Gliner calls global macro ``often considered the most flexible and opportunistic hedge fund strategy'', because of the range of products and markets it trades, and gives its aim as preserving capital ``using stringent risk management to limit drawdowns'' \citep[chap.~1, opening, PDF p.~29]{s7-gliner}. He distinguishes four approaches (sec.\ ``Types of Global Macro Strategies'', PDF pp.~30--33):
\begin{itemize}
\item \emph{discretionary}: the trader's experience and top-down analysis of industries, sectors, countries and the macroeconomy, expressed through directional and relative-value trades (pp.~30--31);
\item \emph{systematic}: model-driven positions without discretion; over several years or more ``systematic funds can produce more consistent returns than discretionary strategies; however, in periods of high volatility, they tend to underperform discretionary macro, as they did in 2008''; trading many liquid markets lets them scale (pp.~31--32). He advises allocating to both, because discretionary macro ``is negatively correlated during periods of stress'' while systematic macro offers predictable, back-tested drawdowns (p.~32);
\item \emph{high-frequency}: holding periods from milliseconds to a few hours, not as large or scalable (p.~32);
\item \emph{CTAs}: mostly automated trend followers, whose position sizing ``originated with the Turtle Traders''; they perform well over several years but suffer large drawdowns (p.~33).
\end{itemize}
\emph{Trades.} A directional trade bets on an asset's direction (long copper on a bullish commodity view); a relative-value trade pairs assets, as in short Italian BTPs against long German Bobls in the euro crisis (p.~31). Thematic trades are long-horizon views, which Gliner sizes smaller and treats separately (chap.~2, sec.\ ``Thematic Trade'', PDF p.~60). \emph{Product groups.} Gliner's building blocks are fixed income, FX, equities and commodities (chap.~4, sec.\ ``The Four Product Groups'', PDF pp.~92--97), which he calls among the most liquid markets; he calls fixed income ``the most important'' of them (p.~95).

\qffield{Notation} \emph{Discretionary} and \emph{systematic} describe how positions are decided, not what is traded; \emph{CTA} is a U.S.\ regulatory category (commodity trading advisor) whose members mostly follow trends.

\qffield{Conditions} The descriptions are Gliner's, as of 2013--2014, and are characterisations, not tested results: the 2008 comparison is a statement without data in the passage, and ``negatively correlated during periods of stress'' is not quantified there (the record is in \hyperref[ev:global-macro-crisis-performance]{global macro in crises}).

\qffield{Practice} \emph{Expressing one view several ways}: a growth view through duration, equity beta, pro-cyclical currencies and cyclical commodities; an inflation view through commodities, breakevens and TIPS; a crisis view through bonds, safe-haven currencies, volatility and gold; a policy-divergence view through rate differentials and curve spreads. This list is an interpretive synthesis, consistent in spirit with Gliner's chapter~4 on relationships between assets (pp.~98--102); the asset responses it assumes are regime-dependent (\hyperref[def:macro-regime-asset-map]{regime map}).

\emph{Liquidity.} The markets Gliner names (G10 and major EM currencies, government bond futures, equity index futures, major commodity futures) are why he calls macro ``one of the most, if not the most, liquid strategies'' (chap.~1, PDF p.~36). Liquidity does not guarantee scale without market impact, fast exit in stress or better redemption terms: impact grows with size even in liquid futures (\hyperref[ev:square-root-impact]{square-root law}), and liquidity can vanish in stress.

\qffield{Pitfalls} \emph{Unsourced explanations.} Explaining the strategy's existence by policy moving in jumps while markets are continuous, by sticky institutional capital or by under-arbitraged cross-asset links, and describing macro as exploiting structural inefficiency rather than informational advantage, or as trading economic regimes rather than securities, are hypotheses, not findings of Gliner or the other sources. \emph{Sizing thematic trades.} Sizing thematic trades smaller because they take years and may carry negatively is a discretionary choice; Gliner discusses them as a separate category with their own sizing (p.~60), and the claim is advice, not a result.

\qffield{See also} \hyperref[ev:global-macro-crisis-performance]{global macro in crises}; \hyperref[meth:macro-position-sizing-risk]{macro position sizing and risk}; \hyperref[meth:macro-practitioner-frameworks]{practitioner frameworks}; \hyperref[def:macro-regime-asset-map]{regime map}; \hyperref[def:fixed-income-universe]{fixed-income universe}.

\qffield{Further reading} Carry across asset classes \citep{alpha-carry} and the section's other readings on trend and value test what Gliner describes.

\end{qfdefinition}

\begin{qfevidence}[Global macro's return record and its behaviour in crises: Gliner's index comparisons, conditional correlation with equities, and what the cited figures do and do not show]\label{ev:global-macro-crisis-performance}
\qffield{Statement} Gliner reports that the Barclays Global Macro Index ``has achieved annualized returns of 10 percent from 2002 to 2012 compared to the S\&P 500, which has been 2 percent'' with lower annualised volatility, and infers ``a higher Sharpe ratio'' than other hedge fund strategies; Figure~1.1 plots the Dow Jones Credit Suisse Global Macro Hedge Fund Index against the S\&P 500 from January 1995 to September 2013 \citep[chap.~1, sec.\ ``Return Profile and Allocations'', PDF p.~34]{s7-gliner}. He states that global macro ``has shown a low correlation to S\&P 500 returns, particularly in periods of market stress'', with ``a negative correlation during market collapses'' because many macro traders go short in bear markets (Figure~1.2, pp.~34--35), and that its return profile and liquidity make it ``the most popular hedge fund allocation by pension funds'' (Figure~1.3, Barclays Prime Services, 2009--2012, pp.~36--38).

\qffield{Notation} Hedge fund indices are averages of reporting funds' net returns; their Sharpe ratios use reported monthly returns.

\qffield{Conditions} These are index statistics over one window, quoted without standard errors, and subject to the usual hedge-fund-index biases (self-reporting, survivorship and backfill), which Gliner does not discuss. The S\&P 500's 2 percent over 2002--2012 depends on whether dividends are included, which the passage does not say. The higher Sharpe ratio than other hedge fund strategies is asserted, not shown.

\qffield{Practice} Macro's record is often summarised as lower volatility than equities, comparable or higher long-run returns, low or negative correlation in crises, and a higher Sharpe ratio, with crises described as times when equity managers are trapped long, credit liquidity vanishes and correlations converge toward 1, while macro funds can cut gross exposure, short directly, and trade rates, FX or volatility. Those mechanisms are plausible but are an interpretive generalisation; Gliner's evidence is the index comparison and the conditional-correlation chart above, and his chapter~4 shows correlations rising around Lehman (Figure~4.5, PDF pp.~100--101).

\qffield{Pitfalls} Do not extend a 2002--2012 or 1995--2013 record to later years: the guide's section readings and later index data are needed. \emph{Own bounded test, 2022--2026}. The protocol predates the data acquisition. The data choices below were fixed before computing. It uses the Barclay Global Macro Index's monthly net returns as displayed on BarclayHedge's public page on 17 October 2026, excluding the months it marks as estimated, against French's Mkt-RF, January 2022 to July 2026 (55 months).
\begin{itemize}
\item \emph{All months.} The index earned 0.68 percent a month at 5.6 percent annual volatility, an annualised Sharpe ratio on excess returns of 0.76 (monthly 0.22, Lo standard error 0.14).
\item \emph{Down months} (24): it averaged $-0.39$ percent ($t=-1.5$), with correlation $+0.29$ (95 percent interval $-0.13$ to $0.62$) and beta $+0.13$.
\item \emph{Up months} (31): it averaged $+1.51$ percent.
\item \emph{The ten worst benchmark months}, which averaged $-6.5$ percent: the index averaged $-0.79$ percent.
\end{itemize}
By the pre-set rule (down-month mean positive with $t>2$), crisis protection is not supported in this window. The low betas say the index fell much less than equities, not that it rose.

\emph{The protocol's window, January 2014 to July 2026}. The public page's Excel download covers the whole history from 1997. Run under the protocol and the data choices above, with both date bounds enforced and the two months marked as estimated (August and September 2026) excluded, it gives for 151 months:
\begin{itemize}
\item \emph{All months.} The index earned 0.46 percent a month at 5.1 percent annual volatility, an annualised Sharpe ratio on excess returns of 0.74 (monthly 0.21, Lo standard error 0.08).
\item \emph{Down months} (53): it averaged $-0.52$ percent ($t=-3.0$), with correlation $+0.42$ (95 percent interval $+0.17$ to $+0.62$) and beta $+0.17$ ($+0.07$ to $+0.28$).
\item \emph{Up months} (98): it averaged $+1.00$ percent.
\item \emph{The ten worst benchmark months}, which averaged $-8.6$ percent: the index averaged $-1.22$ percent.
\end{itemize}
The verdict is unchanged, crisis protection not supported, but the longer window strengthens it. The down-month loss is now distinguishable from zero, and the positive co-movement is significant, where the 55-month interval included zero. The 2022--2026 figures above are the recent-period comparison; the workbook reproduces them exactly.

\emph{Limits of the index.} BarclayHedge's page says the index ``is recalculated and updated real-time on this page as soon as the monthly returns for the underlying funds are recorded''. Only funds that report net returns are included, and the number of funds ``will continue to increase until all funds categorized within the sector have reported''. So recent months are provisional, and earlier values can change as late reporters arrive. The page does not state the weighting, how newly added funds' histories are treated (backfill) or how closed funds are kept (survivorship), and no methodology document is public. The test also gives equal treatment to all down months and uses one index of self-reporting funds.

Negative correlation in crises is a conditional statistic on a few episodes, sensitive to how crises are chosen. Gliner's ``most popular hedge fund allocation by pension funds'' is his reading of one prime broker's survey for 2009--2012, not a market-wide statistic.

\qffield{See also} \hyperref[def:global-macro-strategy]{global macro}; \hyperref[meth:macro-position-sizing-risk]{macro position sizing and risk}; \hyperref[meth:relative-performance-evaluation]{relative performance evaluation}.

\qffield{Further reading} Crisis performance of trend and carry after 2013, for example \citep{wide-harvey-worst-times}.

\end{qfevidence}

\begin{qfmethod}[Macro position sizing and risk limits: Turtle-style volatility units, risk/reward targets, gap risk, VaR utilisation tied to year-to-date P\&L, stress tests, performance ratios and portfolio correlation, as Gliner describes them]\label{meth:macro-position-sizing-risk}
\qffield{Statement} Gliner sizes positions by volatility, following the Turtle Traders: $N$ is ``the 20-day exponential moving average of the True Range'', with true range $\max(H-L,\,H-C_{-1},\,C_{-1}-L)$, and position size is set so that a move of $N$ (or $2N$, for the stop) costs a fixed fraction of capital \citep[chap.~2, secs.\ ``Unit Size'' and ``Volatility Adjusting Position Size'', PDF pp.~50--53]{s7-gliner}. He asks for at least 1:3 risk to reward, for example risking 25~bp to make 75~bp, with targets set before entry and reasons logged; a trader right 40 percent of the time then earns $0.4\times75-0.6\times25=15$~bp a trade (sec.\ ``Risk/Reward'', pp.~54--55). Gap risk, a price change with no trading in between, is in his account one of the biggest risks in sizing (sec.\ ``Gap Risk'', p.~57, with Cyprus, 16 March 2013).

\emph{Risk utilisation} (pp.~66--68): with a maximum 95 percent VaR of 5 percent, ``if you have 50 percent conviction, that doesn't mean you should run 50 percent of your risk utilization''; a limit predetermined at the start of the year rises to 6 percent if the portfolio is up 5 percent and 7 percent if up 10 percent, falls to 4 percent if down 2.5 percent and 3 percent if down 5 percent, and the year starts again at 0 percent return.

\emph{Stress tests} (p.~69): sort the last 500 days' P\&L, average the five worst, and in calm periods also replay extreme days (his Table~2.1 includes the October 1987 crash, the 2011 Japanese earthquake and the U.S.\ downgrade).

\qffield{Notation} $H,L,C_{-1}$ the day's high and low and the previous close; $N$ in price units; VaR as a percentage of capital at 95 percent over one day (``the worst one out of 20 days'', p.~66).

\qffield{Conditions} $N$ is an average range, not a standard deviation. Gliner's statement that a $1N$ stop is hit with about 16 percent probability and a $2N$ stop with about 2.3 percent (pp.~52--53) treats $N$ as a one-day standard deviation; for Brownian daily moves the mean true range is about $2\sqrt{2/\pi}\approx1.6$ standard deviations, so a close-to-close move beyond $2N$ is far rarer than 2.3 percent under normality, and fat tails and gaps change both numbers. The 1:3 rule's break-even hit rate is 25 percent before costs (section~2); it says nothing about whether the targets are reachable. A 5 percent one-day VaR is Gliner's illustration and implies an annualised volatility near 48 percent at normality, far above typical macro fund risk.

\qffield{Practice} Gliner's metrics (sec.\ ``Sharpe Ratio'' to ``Value at Risk'', pp.~62--66): Sharpe ratio, which he says ``does assume a normal distribution''; Sortino ratio, with downside deviation; drawdowns; and VaR. The average of the 5 worst of 500 days is an empirical 1 percent expected shortfall, not a ``three standard deviation'' loss as he puts it (p.~69); under normality it is about 2.7 standard deviations. \emph{Correlation.} Gliner warns that positions correlated with one another can be stopped out together (sec.\ ``Correlation'', p.~56) and that correlations rise in stress (chap.~4, pp.~100--101). His illustration that perfectly correlated assets move by equal amounts (p.~100) confuses correlation with beta: $b=2a$ has correlation one and beta two (checks, section~3).

\qffield{Pitfalls} \emph{Additional scenarios.} LTCM (1998), the 2015 renminbi devaluation and the 2020 COVID crash are reasonable stress scenarios but are not on Gliner's list. Gliner says the Sharpe ratio assumes a normal distribution, but the ratio itself does not, though reading it as a probability statement does. VaR is a loss exceeded on about 1 day in 20 (or 100), not an expected loss. That position sizing matters more than knowing what to buy or sell is advice, not a result. \emph{Raising risk after gains} is a P\&L-contingent rule Gliner links to Thorp; whether it improves risk-adjusted returns is untested here.

\qffield{See also} \hyperref[def:global-macro-strategy]{global macro}; \hyperref[meth:relative-performance-evaluation]{relative performance evaluation}; \hyperref[meth:macro-practitioner-frameworks]{practitioner frameworks}; \hyperref[def:macro-regime-asset-map]{regime map} (correlations by regime).

\qffield{Further reading} Volatility targeting and its out-of-sample evidence are in \hyperref[ev:volatility-managed-portfolios]{volatility-managed portfolios}.

\end{qfmethod}

\begin{qfmethod}[Practitioner frameworks from Drobny's interviews: risk-premia allocation, long optionality, trade construction, front-end rates, commodity microeconomics, equity concentration, the anonymous currency manager, emerging-market crises, flows and contrarian signals]\label{meth:macro-practitioner-frameworks}
\qffield{Statement} These are practitioners' opinions recorded in interviews around 2005, not tested rules; each is attributed to the interviewee who said it in \citep[Wiley 2006 printing, printed pages]{s7-drobny}:
\begin{itemize}
\item \emph{Leitner}: start ``with 20 percent in each'' of five asset classes (equities, fixed income, currencies, commodities, real estate) to earn risk premia (p.~60); fixed income levered at the front end, where ``you could be five times levered in the two-year'' for the risk of a 10-year bond (p.~62); commodities by rolling the second contract, said to pay 6 percent real a year (p.~62); an absolute-return bucket for ``great, out-of-the-box ideas we come across about twice a year'', with up to 10 percent of the fund on one idea (p.~63); look for disconfirming evidence (p.~50); not being invested forgoes risk premia, and fighting them requires being ``doubly right'' (p.~54).
\item \emph{Siva-Jothy}: never short gamma, carry being synthetically short gamma (pp.~78--79); ``I always structure my portfolio to be 50 to 60 percent options'', with three-to-six-month options and overnight options, without saying whether the share is of notional, premium or risk; ``three or four major macro opportunities each year''; ``Overtrading is as bad as running losing positions for too long'' (p.~85).
\item \emph{Andres Drobny}: think about positioning over ``the last two, three, or five years'', and a rising current-account deficit means ``a long building up'' (p.~109); agrees with the interviewer's ``three or four good trades a year'' style, looking for the few 20--30 percent events with about 80 percent of the money in carry or lower-risk positions (p.~125).

\emph{Stops} (pp.~131--132): after a stop on a single trade you may get back in, ``but with a very tight stop'', counting the loss already taken, so the risk/reward, which ``should at least be 4 to 1 in any given trade'', is worse the second time; after hitting the 1 ``a couple of times'' a 4:1 trade has become 4:3 and you should move on. Asked about a cooling-off period, he answers ``On a specific trade, no''; only ``On a portfolio level, if you get stopped out of all positions due to a bad run'', after big losses, should managers ``get away from the markets for about a month''; he also recommends three to four weeks away once a year. \emph{Wadhwani}: positioning matters more than the macro news (p.~163); never ignore a stop, and take a 24-hour break before re-entering (p.~177).
\item \emph{Porter}: ``time horizon arbitrage'' against short-term investors (p.~135); maximise current income per unit of VaR at the front end; expertise in one-year, one-year forwards, ``the most volatile part of the yield curve''; ``most interest rate cycles are between 16 and 24 months'' (p.~139; p.~148); a steepener means ``you basically think the bonds are going to rally'' and a flattener that ``the market is going to sell off'' (pp.~147--148); central banks and magazine covers as counter-indicators (p.~152).
\item \emph{Rogers}: a Fortune cover or front-page story as a sign a move is over (pp.~217, 230).
\item \emph{Anderson}: invest from microeconomics and fundamentals, not price (pp.~249, 256); stress the thesis against adverse macro (``when we're long, we assume very low growth'', pp.~256--257); ``margins revert to the mean'', not prices (p.~258); size by volatility (p.~259); only low-cost producers, since ``you can't buy a high-cost asset cheap enough'' (p.~261); research on the ground (pp.~261--262); managing off a VaR number is ``incredibly dangerous'' for a concentrated book (p.~250).
\item \emph{Bessent}: ``50 percent of a stock's move is the overall market, 30 percent is the industry group, and then maybe 20 percent'' stock-specific (p.~280); 8 to 14 large, unequally weighted positions (p.~283).
\item \emph{Dimitrijević}: country limits of 10 percent in the macro fund and 15 percent in the emerging-market fund (p.~293); the clean-sheet question, would we buy at the current price (p.~294); hedge a rich asset against an attractive one in the same country, short Brazilian debt against Brazilian equities (p.~295); stops ``work in continuous markets'' and are dangerous in distressed or value positions, are kept internal because stops given to counterparties ``invariably get hit'', and are followed by a one-month cool-off (pp.~295--296); the best EM opportunities come after panics (p.~296).
\item \emph{Yra Harris}: the ``matrix'' of global money flows seeking the highest risk-adjusted return; ``money \dots\ craves stability''; fear and greed (p.~207); the Brazilian real devaluation of 1999 and selling soybeans (p.~208); Dennis's Slower Fool Theory (pp.~204--205); never the ``prayer trade'' back to break-even; ``clear the books, and start over'' (pp.~209--210); best when focused on three or four things (p.~214); ``80 percent of my day is being wrong'', probing ``like a reconnaissance mission'' (p.~215).
\item \emph{The anonymous currency specialist}: FX as the tail of the credit curve (p.~327); three blocs, Europe, Asia and Latin America (p.~328); three or four macro calls (pp.~328--329); a gestalt shift (p.~329); no relative value, from a fear of negative gamma (p.~330); bullish means long or flat (p.~331); long gamma and vega, leverage a function of correct analysis, ``75 percent of price movement'' in ``10 to 15 percent of the time'' (p.~332); never a forced seller (p.~333); a five-month horizon to judge a view (p.~334); the ``frozen ice'' theory of crowded trades (p.~336).
\end{itemize}

\qffield{Notation} Page numbers are printed pages of the 2006 printing (copyright 2006, ISBN 978-0-471-79447-9); the catalogue entry describes a revised 2013 edition, which was not inspected.

\qffield{Conditions} Opinions of successful survivors, selected by the author, with no record of how the rules performed; several conflict (Leitner's long risk premia against the anonymous manager's fear of carry as negative gamma; Andres Drobny's 4:1 against Gliner's 1:3).

\qffield{Practice} Attributions in the interviews: the ``75\% of price moves'', long-gamma leverage, forced-seller and ``inevitable'' timing points are the anonymous manager's (pp.~331--333, 336); 4:1 and the stop rules are Andres Drobny's, and his cooling-off applies to a portfolio-wide stop-out after large losses, not to a single trade; ``3--4 major opportunities annually'' and ``overtrading as damaging as running losers'' are Siva-Jothy's (p.~85), while Andres Drobny only accepts the interviewer's three-or-four-trades framing (p.~125); Wadhwani's own re-entry rule is a 24-hour break (p.~177); time-horizon arbitrage is Porter's; contrarian front-page signals are Rogers's and Porter's.

\qffield{Pitfalls} \emph{Reading the interviews.} Porter says the curve ``is 90 percent explained by the direction of interest rates'' (pp.~147--148), which does not make a steepener a 90 percent bullish bond trade. The anonymous manager's three blocs are Europe, Asia and Latin America. Leitner's FX carry and volatility-premium trades and leveraged short-duration fixed income are examples, not rules, and his equal 20 percent risk allocation is a starting point, adjusted afterwards. The contrarian observations come from Porter, Rogers and the anonymous manager.

\qffield{See also} \hyperref[def:global-macro-strategy]{global macro}; \hyperref[meth:macro-position-sizing-risk]{macro position sizing and risk}; \hyperref[meth:curve-trades-carry-rolldown]{curve trades} (steepeners and direction); \hyperref[mod:fx-carry-trade]{FX carry}.

\qffield{Further reading} Drobny's later collections and the section's systematic evidence.

\end{qfmethod}
